% TOP_PROBLEM: {"id": 3293, "title": "¿Are critical k-forests tight?", "category_id": 3, "category": {"id": 3, "name": "graph_theory", "display_name": "Graph Theory", "description": "Problems involving graphs, networks, and their properties.", "slug": "graph-theory", "order_index": 3, "created_at": "2026-07-31T15:26:25.671Z"}, "status": "open"}
% FIRST_POSTED: null
% ENTRY_KIND: statement_only
% Imported from: batch03.tex; UnsolvedMath ID 3293
\documentclass[11pt]{article}
\usepackage[margin=25mm]{geometry}
\usepackage{fontspec}
\setmainfont{Latin Modern Roman}
\usepackage{amsmath,amssymb,mathtools}
\usepackage{unicode-math}
\setmathfont{Latin Modern Math}
\usepackage{xurl,hyperref}
\hypersetup{colorlinks=true,urlcolor=blue}
\setcounter{secnumdepth}{0}
\setlength{\parindent}{0pt}
\setlength{\parskip}{5pt}
\setlength{\emergencystretch}{3em}
\newcommand{\sref}[2]{\hyperlink{source:#1:#2}{[S#2]}}
\newcommand{\cataloguescope}[1]{\paragraph{Recorded target.}#1}
\begin{document}
% UnsolvedMath ID 3293; source catalogue code OPG-547
\setcounter{section}{3292}
\section{¿Are critical k-forests tight?}\label{3293}
{\small\sffamily UnsolvedMath ID 3293 · Graph Theory}
\subsection{Definitions and mathematical statement}

\paragraph{Shared notation.}$\mathbb N=\{1,2,\ldots\}$, and $\mathbb Z,\mathbb Q,\mathbb R,\mathbb C$ denote the integers, rationals, reals and complex numbers. Any other conventions are specified locally.


Fix $k\geq2$ and a finite set $V$ with $|V|\geq k$. A $k$-uniform hypergraph is $H=(V,E)$ with $E\subseteq\binom Vk$. A $k$-colouring means a surjection $c:V\to[k]=\{1,\ldots,k\}$. An edge is rainbow when $c(e)=[k]$. The hypergraph is a $k$-forest if each $e\in E$ is the unique rainbow edge for some $k$-colouring. It is critical (saturated) if adding any missing $k$-element edge on the same vertex set destroys this property. It is tight if every surjective $k$-colouring has a rainbow edge; a tight $k$-forest is a $k$-tree.
\paragraph{Statement recorded in the supplied source.}
Must every critical $k$-forest be tight, and hence be a $k$-tree? This is the historical conjecture; the cited primary paper gives a saturated $3$-forest on six vertices which is not tight. \sref{3293}{2} \sref{3293}{3}
\subsection{Short English statement}
Does being inclusion-maximal among these hypergraph forests force the corresponding connectivity property? The historical assertion has a six-vertex counterexample.
\subsection{Sources}
\begin{description}
\item[\hypertarget{source:3293:1}{S1}] ULAM.AI and the UnsolvedMath Contributors, UnsolvedMath: A Curated Collection of Open Mathematics Problems, record 3293 (OPG-547). CC BY 4.0. \url{https://huggingface.co/datasets/ulamai/UnsolvedMath}.
\item[\hypertarget{source:3293:2}{S2}] Open Problem Garden, ¿Are critical k-forests tight?, node 547, original problem page. Original question and full discussion. \url{https://www.openproblemgarden.org/op/are_critical_k_forests_tight}.
\item[\hypertarget{source:3293:3}{S3}] Heidi Gebauer, Anna Gundert, Robin A. Moser and Yoshio Okamoto, Not All Saturated 3-Forests Are Tight (2011), Section 2 and Theorem 2.2. \url{https://arxiv.org/pdf/1109.3390}.
\end{description}
\label{end:3293}
\end{document}
